\section{Reproducibility and integrity conditions}
\label{app:integrity}

The frozen ontology contains seven leaves: built structure, transport surface, woody vegetation, herbaceous vegetation, cropland, bare surface, and surface water. Herbaceous vegetation is the frozen output-index name; its OEM target relation is the broader rangeland vegetation proxy, including bushes/shrubs, while the Semantic Drone relation is grass. These output aliases do not define a fully reconciled botanical taxonomy. Cropland denotes the declared agricultural-field or managed-cover relation. Under the frozen crosswalk, LoveDA building, road, forest, agriculture, barren, and water provide exact operational support for six non-herbaceous leaves. FLAIR raw ID 10 is the only FLAIR label admitted by B1. OpenEarthMap rangeland and Semantic Drone grass are declared target relations. OpenEarthMap developed space remains a coarser ancestor, not building or road support.

\begin{table}[h]
\centering
\caption{Complete frozen prompt set used by every reported arm.}
\label{tab:prompts}
\begin{tabular}{ll}
\toprule
Common leaf & Prompt \\
\midrule
Built structure & a satellite image of a building \\
Transport surface & a satellite image of a road \\
Woody vegetation & a satellite image of forest or trees \\
Herbaceous vegetation & a satellite image of grassland or rangeland \\
Cropland & a satellite image of agricultural cropland \\
Bare surface & a satellite image of bare soil or barren ground \\
Surface water & a satellite image of water \\
\bottomrule
\end{tabular}
\end{table}

Each raster is partitioned at stride 512 into non-overlapping \(512\times512\) RGB tiles. Only the right and bottom edge tiles are zero-padded. CLIP patch features are scored, logits are bilinearly interpolated to \(512\times512\) with the align-corners parameter set to false, and the leaf argmax is copied over the unpadded tile extent. No overlap, blending, test-time augmentation, or post-processing is used.

For B0-half-batch2, duplicate copies of the same deterministic crop have the same individual loss. With mean loss, a frozen backbone, no train-mode batch-dependent state, no sample-dependent random transformation at the adapter, and identical optimizer state, the mean gradient equals the one-crop gradient at every update. The absolute mean-mIoU contrast is below \(4\times10^{-9}\), numerically zero at the reported precision; it is therefore a constructive pipeline-integrity check, not an empirical estimate of a generic batch-size effect.

\begin{table}[h]
\centering
\caption{Duplicated-crop pipeline-integrity check. B0-half-batch2 duplicates the deterministic LoveDA crop in the second update slot.}
\label{tab:batch_integrity}
\begin{tabular}{@{}p{0.28\linewidth}p{0.30\linewidth}p{0.27\linewidth}@{}}
\toprule
Contrast & Metric & Audited result \\
\midrule
B0-half-batch2 minus B0-half & absolute mean mIoU difference & below \(4\times10^{-9}\) \\
\bottomrule
\end{tabular}
\end{table}

Table~\ref{tab:all_leaf_contrasts} provides the complete descriptive, unadjusted intervals for the seven OpenEarthMap leaves.

\begin{table}[h]
\centering
\caption{Unadjusted descriptive paired city-bootstrap intervals for all OpenEarthMap leaf contrasts. The primary target leaf is herbaceous/rangeland.}
\label{tab:all_leaf_contrasts}
\begin{tabular}{lrrr}
\toprule
Leaf & B1 minus B0 IoU & Lower & Upper \\
\midrule
Built structure & -0.01958 & -0.04881 & +0.00555 \\
Transport surface & -0.1482 & -0.1845 & -0.1134 \\
Woody vegetation & -0.1554 & -0.1880 & -0.1225 \\
Rangeland proxy & +0.3118 & +0.2821 & +0.3399 \\
Cropland & -0.2005 & -0.2497 & -0.1498 \\
Bare surface & -0.05080 & -0.1094 & -0.00103 \\
Surface water & -0.03249 & -0.07099 & +0.00676 \\
\bottomrule
\end{tabular}
\end{table}
